\documentclass[11pt]{article}
\usepackage[margin=1in]{geometry}
\usepackage{amsmath,amssymb,amsthm,booktabs}
\usepackage[hidelinks]{hyperref}
\usepackage{xurl}
\setlength{\emergencystretch}{2em}
\newtheorem{theorem}{Theorem}
\newtheorem{lemma}[theorem]{Lemma}
\newcommand{\mr}{\operatorname{mr}}
\title{Twenty-four sharp bounds from simple $A_8$ quotients}
\author{Written proof and exact finite certificates}
\date{17 November 2026}
\begin{document}\maketitle

Twenty-four knot groups admit epimorphisms onto $A_8$ taking every
meridian to a double transposition. This conjugacy class requires four
generators although $A_8$ itself requires only two. The resulting
meridional lower bound is elementary: three double transpositions
cannot act transitively on eight points. Literal permutation assignments
and four-seed Wirtinger traces give independent proofs that the knots
have meridional rank and bridge number four.

These are the same twenty-four named knots previously certified with
ternary affine quotients. The present target has order $20160$, smaller
than the order $122472$ of those targets. The present lower proofs use
neither the older affine certificates nor any branched-cover homology.
No additional named equality is claimed.

\section{The finite group and its meridian class}

Let $C$ be the double-transposition class in the natural alternating
group $A_8$ on $\{0,\ldots,7\}$. The generating rank of $C$ is the
least number of its members generating $A_8$.

\begin{lemma}\label{lem:transitive}
If $r$ double transpositions generate a transitive permutation group
on eight points, then $r\ge4$.
\end{lemma}
\begin{proof}
Make a graph on the eight points, putting in the two transposition
edges from each generator. The connected components are exactly the
generator orbits. Transitivity therefore requires a connected graph,
which has at least seven distinct edges. There are at most $2r$
edges, so $7\le2r$.
\end{proof}

\begin{theorem}\label{thm:group}
The simple group $A_8$ has order $20160$ and ordinary generating
rank two. Its class $C$ has $210$ elements and generating rank exactly
four.
\end{theorem}
\begin{proof}
The order is $8!/2=20160$. There are $\binom84\cdot3=210$ double
transpositions. They form one $A_8$ class: they form one $S_8$ class,
and an odd conjugator can be made even by multiplying by one of the
two transposition factors centralizing the source element.

For an attaining quartet use
\[
\begin{aligned}
 x_0&=(0\,1)(2\,3),&x_1&=(1\,4)(2\,3),\\
 x_2&=(0\,5)(1\,6),&x_3&=(1\,2)(6\,7).
\end{aligned}
\]
The exact finite enumeration described below gives
$|\langle x_0,x_1,x_2,x_3\rangle|=20160$. All four generators are even,
so they generate $A_8$. Lemma~\ref{lem:transitive} excludes three
$C$-generators and proves the class rank is four.

The two even permutations
\[
             a=(0\,1\,2\,3\,4\,5\,6),\qquad b=(0\,1\,7)
\]
also generate a group of order $20160$, as verified by the same exact
enumeration. Thus the ordinary rank is at most two; noncommutativity
excludes rank one. For completeness, conjugating the full enumerated group by the
ordinary pair and their inverses partitions it into classes of sizes
\[
1,105,112,210,1120,1260,1344,1344,1344,1680,2520,2880,2880,3360.
\]
An exact subset-sum check finds no union of classes containing the
identity with order a proper nontrivial divisor of $20160$.
Every normal subgroup is such a union and satisfies Lagrange's
order condition, proving simplicity. Simplicity is not needed for
the meridional lower bound.

For clarity, each enumeration begins at the identity, repeatedly
right-multiplies every newly found permutation by each generator and
its inverse, and stops when no new permutation appears. The retained
set is then precisely the generated subgroup. The permutations are
literal eight-entry tuples, so all products, equality tests and the
final count are exact. No group-recognition or classification software
is used.
\end{proof}

\section{Twenty-four named epimorphisms}

\begin{theorem}\label{thm:sharp}
Every knot in
\[
\begin{aligned}
\mathcal K=\{&13a_{373},13a_{375},13a_{391},13a_{392},\\
&13a_{3264},13a_{3266},13a_{3275},13a_{3276},\\
&13n_{1005},13n_{1006},13n_{1007},13n_{1011},13n_{1012},13n_{1013},\\
&13n_{1043},13n_{1044},13n_{1045},13n_{1046},13n_{1047},13n_{1048},\\
&13n_{3842},13n_{3844},13n_{3855},13n_{3856}\}
\end{aligned}
\]
admits an epimorphism $\pi_1(S^3\setminus K)\twoheadrightarrow A_8$
taking meridians into $C$. For every such knot,
\[
                         \mr(K)=b(K)=4.
\]
\end{theorem}
\begin{proof}
The retained input supplies a literal thirteen-crossing PD, thirteen
permutations, four seed arcs and a complete nine-step coloring trace
for each name. Number PD slots from zero. Join the edge labels in
slots one and three through every overpass, sort the resulting
connected components lexicographically, and number these Wirtinger
arcs $0,\ldots,12$. A PD row $(a,b,c,d)$ has over-arc $[b]=[d]$
and under-arcs $[a],[c]$.

Every retained permutation $p_i$ is a double transposition.
For every reconstructed crossing $(o,j,k)$ the checker verifies
\[
                          p_op_jp_o=p_k.
\]
Since $p_o=p_o^{-1}$, this satisfies either signed Wirtinger convention.
The assignments therefore define knot-group homomorphisms taking all
meridians into $C$. The four seed images generate a group of order
$20160$ in every case, proving surjectivity onto $A_8$.
For $13a_{373}$ the seed images are the displayed attaining quartet.

Any meridional generating family maps to a $C$-generating family of
$A_8$: meridians are conjugate up to inversion, and $C=C^{-1}$.
Theorem~\ref{thm:group}, or the transitivity argument alone, gives
$\mr(K)\ge4$.

Each retained trace starts with four distinct seed arcs and reaches
all thirteen arcs, at each step using a known over-arc and a known
under-arc to color the other under-arc. Thus the diagram has Wirtinger
number at most four. The published bridge--Wirtinger equality
\cite[Definition~1.2 and Theorem~1.3]{bkvv} gives $b(K)\le4$.
The elementary bridge presentation yields $\mr(K)\le b(K)$, proving
both equalities. The bridge upper bound uses this exact coloring
trace and the published theorem; meridional generation alone would
not give it.

The saved named-census PDs are tied to the literal diagrams by exact
dart bijections preserving edge gluing, vertex rotation and over/under
information, with a possible global mirror or orientation reversal.
These operations preserve both ranks and the involutory meridian
class. The independent census names remain an explicit trusted import;
they are not inferred from the finite-group calculation.
\end{proof}

\section{Exact certificates and scope}

The necessary input is \path{a8_simple_quotient_inputs.json}. Run
\path{verify_a8_simple_quotients.py} with the standard Python library.
It checks the sphere embeddings and one-component condition, all
named dart maps, all $312$ literal crossing relations, all meridian
classes, all four-seed traces, all image orders, the $210$-element
conjugacy orbit, the ordinary generating pair and the attaining
quartet. Corrupted permutations, traces, named maps, ordinary pairs
and class quartets are rejected.

The optional \texttt{--census} mode rereads all twenty-four named PDs
and DT codes from the pinned \texttt{snappy\_15\_knots}~1.2.1 database
using Spherogram~2.4.1, checking the exact database hash and size.
The default replay checks the retained records and their exact map
certificates. The published primary bridge--Wirtinger theorem is
retained with a source hash and is imported rather than reproved.

Only positive epimorphisms are asserted here. A complete negative
classification on other knots is not part of this result. Ordinary
rank two concerns the finite target $A_8$, with no corresponding
ordinary-rank claim for the knot groups. The computation is an internal
exact replay, not a formal verification or an external specialist
review. No worldwide-priority claim is made.

\section{A representative literal assignment}

For $13a_{373}$ the four retained seed arcs are $0,1,5,10$. Their images
are $x_0,x_1,x_2,x_3$ in that order. The entire arc assignment is
\[
\begin{array}{c|c@{\qquad}c|c}
\text{arc}&\text{permutation}&\text{arc}&\text{permutation}\\\hline
0 & $(0\,1)(2\,3)$ & 7 & $(2\,6)(3\,5)$\\
1 & $(1\,4)(2\,3)$ & 8 & $(0\,6)(3\,7)$\\
2 & $(0\,4)(2\,3)$ & 9 & $(1\,5)(3\,7)$\\
3 & $(0\,4)(5\,6)$ & 10 & $(1\,2)(6\,7)$\\
4 & $(0\,1)(5\,6)$ & 11 & $(0\,2)(5\,7)$\\
5 & $(0\,5)(1\,6)$ & 12 & $(0\,2)(1\,3)$\\
6 & $(2\,5)(3\,6)$ & &\\
\end{array}
\]
The literal PD, with rows in crossing order, is
\[
\begin{gathered}
(3,1,4,26),\ (1,7,2,6),\ (7,3,8,2),\\
(13,4,14,5),\ (5,14,6,15),\ (21,9,22,8),\\
(9,17,10,16),\ (17,11,18,10),\ (11,25,12,24),\\
(19,12,20,13),\ (15,22,16,23),\\
(23,18,24,19),\ (25,21,26,20).
\end{gathered}
\]
Using the arc-numbering rule in the proof, the complete coloring trace
is the following list of triples $(o,a,b)$: the known over-arc $o$ and
known under-arc $a$ color the previously unknown under-arc $b$.
\[
\begin{gathered}
(0,1,2),\ (10,0,12),\ (12,5,6),\\
(6,10,9),\ (9,12,11),\ (2,6,7),\\
(7,2,3),\ (5,9,8),\ (1,3,4).
\end{gathered}
\]
These displayed data give one complete finite certificate without
consulting the machine supplement. The other twenty-three cases are
stored in the same literal format in the input.

\begin{thebibliography}{9}
\bibitem{bkvv}
R. Blair, A. Kjuchukova, R. Velazquez and P. Villanueva,
Wirtinger systems of generators of knot groups,
Commun. Anal. Geom. 28 (2020), 243--262,
Definition~1.2 and Theorem~1.3, printed p.~246.
\url{https://archive.intlpress.com/site/pub/files/_fulltext/journals/cag/2020/0028/0002/CAG-2020-0028-0002-a002.pdf}.
\end{thebibliography}
\end{document}
